\subsection{The conjugate pair and its actual wild descent}

\begin{statement}{E.80}{Proved}{The joint cubic quotient, its integral filtration, and complete descended local fibres}
Put $K=\mathbb Q(\zeta)$, $\zeta=\zeta_3$, and
$k=\mathbb Q_3(\zeta)$. Retain the cubic, tangent and
actual source marking of E.78. In particular
$O=[1:-1:0]$, the source parameter is $s=z/x$ of value
one, and $N$ is the based geometric free pro-three
group of $U=\{x^3+y^3+z^3=0,\ xyz\ne0\}$.
Use conjugate paths for the two maps
\[
 \begin{gathered}
 h=\frac{x+y}{x+\zeta y},\qquad
 \bar h=\frac{x+y}{x+\zeta^2y},\qquad
 v=-\frac{x+y}{z},\\
 \bar h=-\frac{\zeta^2h}{h+\zeta},\qquad
 \frac{1-h}{1-\bar h}=-\zeta\frac h{\bar h},\qquad
 v^3=h\bar h.
 \end{gathered}
 \tag{E80.1}
\]
Both target groups are the same marked free pro-three
group $F=\langle x,y\rangle$ of the three-punctured
line. Write $X(f),Y(f)$ for its abelian coordinates
and $\ell=X-Y$. The same notation denotes the
characters of $C_r=F/\gamma_{r+1}F$ for $1\le r\le3$.

There is the following exact intermediate cover of
schemes over $K$:
\[
 \begin{gathered}
 U\longrightarrow X_0=E\setminus\{yz=0\}
   \longrightarrow S=\mathbb P^1_h\setminus\{0,1,\infty,-\zeta\},\\
 X_0\simeq
 \left\{(h,v):v^3=-\frac{\zeta^2h^2}{h+\zeta}\right\}
       \longrightarrow S.
 \end{gathered}
 \tag{E80.2}
\]
The first arrow fills the three cusps with $x=0$.
The second is connected, cyclic, finite \'{e}tale of
degree three. Let $H=\pi_1^{(3)}(S)$, let
$N_0=\pi_1^{(3)}(X_0)\subset H$, and let
$K_\infty,K_{-\zeta}$ be the closed normal subgroups
generated by the indicated peripheral generators of
$H$. Then $H$ is free of rank three, $N_0$ is its
index-three subgroup of rank seven, and the exact
joint image is
\[
 \begin{gathered}
 \mathcal J=\operatorname{im}(N\longrightarrow F\times F)
   =\{(f,g):\ell(f)=\ell(g),\ X(f)+X(g)\in3\mathbb Z_3\},\\
 \ker(N\longrightarrow\mathcal J)
   =(N\longrightarrow N_0)^{-1}(K_\infty\cap K_{-\zeta}),\\
 K_\infty\cap K_{-\zeta}
      =\overline{[K_\infty,K_{-\zeta}]}\subset N_0.
 \end{gathered}
 \tag{E80.3}
\]
In particular the two arithmetic factors are retained
with their actual joint image. For each $r=1,2,3$ put
\[
 \begin{gathered}
 \mathcal J_r=
 \{(f,g)\in C_r^2:\ell(f)=\ell(g),\ X(f)+X(g)\in3\mathbb Z_3\},\\
 N/\gamma_{r+1}N\twoheadrightarrow\mathcal J_r,\qquad
 \ker=
 \frac{h_*^{-1}(\gamma_{r+1}F)\cap
             \bar h_*^{-1}(\gamma_{r+1}F)}{\gamma_{r+1}N}.
 \end{gathered}
 \tag{E80.4}
\]
These are actual $G_{\mathbb Q_3}$-equivariant quotient
maps with the action below.

Let $3_n,\zeta_n=\zeta_{3^{n+1}},\lambda_n,a_n$ be
the coherent roots in E.78, and use their complex
conjugates for $\bar h$. They satisfy the literal
identities
\[
 \begin{gathered}
 \lambda_n/\bar\lambda_n=-\zeta_n,\qquad
 \lambda_n\bar\lambda_n=3_n,\qquad
 a_n/\bar a_n=-\zeta_n^{-1},\qquad
 a_n\bar a_n=3_n^{-3},\\
 \kappa_a=-\frac{\chi-1}{6}-\frac32\kappa_3(3),\qquad
 \kappa_{\bar a}=\frac{\chi-1}{6}-\frac32\kappa_3(3).
 \end{gathered}
 \tag{E80.5}
\]
The product of the chosen cube roots of $h,\bar h$
is exactly $v$, whose leading term is $s^2/3$.
Thus (E80.2) uses the actual sheet selected by the
source tangent.

Write $c=[x,y]$, $d=[x,c]$, $e=[y,c]$ and
$x^Ay^Bc^nd^re^s$ for the Hall convention of E.78.
Let $T_\mu$ be the polynomial operator (E78.6), now
evaluated at all four values below, and let
$I(x)=x^{-1}, I(y)=y^{-1}$. The actual arithmetic
operators of $\tau,\tau'$ and the order-two
arithmetic generator in the actual relative source
of E.74 and E.77 are respectively
\[
 \begin{aligned}
 \mathcal T(f,g)&=(T_{-5/2}f,T_{-1/2}g),\\
 \mathcal T'(f,g)&=(T_{1/2}f,T_{5/2}g),\\
 \jmath(f,g)&=(Ig,If),\qquad
              \jmath\mathcal T\jmath=\mathcal T'.
 \end{aligned}
 \tag{E80.6}
\]
The operators descend through every map in (E80.4).
Their source generators are the actual E.77
generators after the E.78 Frobenius transport.

Normalize a local cocycle on that relative
order-two subgroup to one. Its complete section
set with coefficients in $\mathcal J_r$ is
\[
 \begin{gathered}
 \{(u_0,a_0)\in\mathcal J_r^2:
       V=u_0\mathcal T(a_0)u_0^{-1}a_0^{-7},\quad
       \jmath(V)=V\}/\mathcal J_r^{\jmath},\\
 w_0:(u_0,a_0)\longmapsto
       (w_0u_0\mathcal T(w_0)^{-1},w_0a_0w_0^{-1}),\\
 \mathcal J_r^{\jmath}
       =\{(w,Iw):\ell(w)=0\}.
 \end{gathered}
 \tag{E80.7}
\]
Restriction to $G_k$ followed by the first target
projection induces an injection
\[
       H^1(\mathbb Q_3,\mathcal J_r)
                  \lhook\joinrel\longrightarrow H^1(k,C_r).
 \tag{E80.8}
\]
Its exact image is given on any E.78 representative
$(u,u',a,a')$ by the four conditions
\[
 \begin{gathered}
 \ell(a)+\ell(a')=0,\qquad X(a)-X(a')\in3\mathbb Z_3,\\
 H_u=\ell(u)+\ell(u')\in3\mathbb Z_3,
       \qquad X(u)-X(u')\in3\mathbb Z_3.
 \end{gathered}
 \tag{E80.9}
\]
Here every representative still satisfies E.78's
actual local source relation. For an explicit inverse,
apply its gauge by $w=x^{H_u/12}$, so the new
horizontal $\ell$-sum is zero, and set
\[
             u_0=(u,Iu'),\qquad a_0=(a,Ia').
 \tag{E80.10}
\]
The exact remaining gauge is $\ker\ell$, as in
(E80.7). Thus (E80.8) includes the full
geometric-conjugacy quotient, and its inverse retains
the actual normalizing element.

The abelian joint lattice has the integral coordinates
\[
 \begin{gathered}
 (L,p,m)\longmapsto
       ((p,p-L),(3m-p,3m-p-L)),\\
 [\jmath]_{L,p,m}=
 \begin{pmatrix}-1&0&0\\0&1&-3\\0&0&-1\end{pmatrix},
       \qquad [\mathcal T]=7I_3.
 \end{gathered}
 \tag{E80.11}
\]
Consequently the complete first-stage coordinates and
the exact next lifting image are
\[
 \begin{gathered}
 H^1(\mathbb Q_3,\mathcal J_1)\simeq
           \mathbb Z_3^5\times\mathbb Z/3,\\
 (L_a,p_a,m_a,L_u,m_u,\ p_u\bmod3),\\
 \operatorname{im}\bigl(H^1(\mathbb Q_3,\mathcal J_2)
                  \to H^1(\mathbb Q_3,\mathcal J_1)\bigr):
       \quad L_a(p_u-1)-L_up_a\equiv0\pmod3.
 \end{gathered}
 \tag{E80.12}
\]
Choose $p_u\in\{0,1,2\}$ for the inverse and use
the fixed gauge with $p=(p_u^{\rm old}-p_u)/6$.
For a particular second-stage lift use the E.78
substitution $n=-D_2/6$, $n'=N=N'=0$ after (E80.10).
Every nonempty fibre at this stage is
$\mathbb Z_3^2\times\mathbb Z/3$, with exactly the
scalar representative and inverse maps of E.78.

At the next stage apply (E78.10)--(E78.11) to the
restricted quadruple. The complete image and fibres
are
\[
 \begin{gathered}
 \operatorname{im}\bigl(H^1(\mathbb Q_3,\mathcal J_3)
                 \to H^1(\mathbb Q_3,\mathcal J_2)\bigr):
       \quad\Delta R,\Delta S\in3\mathbb Z_3,\\
 \text{each nonempty fibre}\simeq
               \mathbb Z_3^4\times(\mathbb Z/9)^2,\\
 a\longmapsto a d^{-\Delta R/336}e^{-\Delta S/336}
                  \quad\text{for a particular lift}.
 \end{gathered}
 \tag{E80.13}
\]
All other representatives and their inverses are
(E78.15), followed by (E80.10). These are
homeomorphisms of compact section fibres.

For example the actual descended second-stage family
\[
 \begin{gathered}
 a=y,\qquad a'=y^{-1},\qquad
 u=x^{5/2}c^N,\qquad u'=x^{-1/2}y^2c^M,\\
 \text{class-three error}=e^{7(5-N-M)}
 \end{gathered}
 \tag{E80.14}
\]
has a third-stage lift exactly for $N+M\equiv2\pmod3$.
A particular lift changes $a$ to
$y e^{(N+M-5)/48}$. Thus the member $N=M=0$ has
empty fibre and $N=2,M=3$ already lifts with its
displayed words. These are sections of the actual
joint quotient over $\mathbb Q_3$; their intrinsic
preimages are governed by the explicit kernels in
(E80.4).

Finally the integral lower central comparison is
as follows. The kernel of
$\mathcal J_2\to\mathcal J_1$ is all of
$\mathbb Z_3^2$ in the two $c$ coordinates, while
\[
 \begin{aligned}
 \gamma_2\mathcal J_2
   &=\{(c^n,c^{n'}):n+n'\in3\mathbb Z_3\},\\
 \gamma_2\mathcal J_3
   &=\{(c^nd^re^s,c^{n'}d^{r'}e^{s'}):
                            n+n'\in3\mathbb Z_3\},\\
 \gamma_3\mathcal J_3
   &=\{(d^re^s,d^{r'}e^{s'}):
          r'-r\in3\mathbb Z_3,\quad
          s+s'-2r\in3\mathbb Z_3\}.
 \end{aligned}
 \tag{E80.15}
\]
In particular the actual comparison maps are
\[
 \begin{gathered}
 0\longrightarrow\mathbb Z/3
   \longrightarrow\mathcal J_2^{\rm ab}
   \longrightarrow\mathcal J_1\longrightarrow0,
       \qquad (n,n')\longmapsto n+n'\bmod3,\\
 1\longrightarrow(\mathbb Z/3)^2
   \longrightarrow\mathcal J_3/\gamma_3\mathcal J_3
   \longrightarrow\mathcal J_2\longrightarrow1,\\
 (r,s,r',s')\longmapsto
           (r'-r,\ s+s'-2r)\bmod3.
 \end{gathered}
 \tag{E80.16}
\]
The two coordinate lattice Smith lists are $(1,3)$
and $(1,1,3,3)$. As compact $\mathbb Z_3$-modules,
$\mathcal J_3^{\rm ab}\simeq\mathbb Z_3^3\oplus\mathbb Z/3$
and
$\gamma_2\mathcal J_3/\gamma_3\mathcal J_3
 \simeq\mathbb Z_3^2\oplus(\mathbb Z/3)^2$.
The image of the intrinsic central map
$\operatorname{Lie}_3(N^{\rm ab})\to
 \ker(\mathcal J_3\to\mathcal J_2)=\mathbb Z_3^4$
is exactly the last lattice in (E80.15), of index
nine. These exact sequences specify how the product
truncation used for (E80.12)--(E80.13) compares with
the actual lower central filtration.

\emph{Manuscript: \S51.16, folios733v--736v,
scans172--178, formulas(371)--(379), Part E(c),(d)
and follow-up(4). The result evaluates the conjugate
joint quotient and its full wild descent; (E80.4)
retains the intrinsic kernel.}
\end{statement}

\begin{proof}
\textbf{The exact schemes and their groups.}
Writing $t=x/y$ gives
$h=(t+1)/(t+\zeta)$ and
$\bar h=(t+1)/(t+\zeta^2)$. Substitution and
$1+\zeta+\zeta^2=0$ prove all three identities
in (E80.1). The inverse to the second line of
(E80.2) is the literal morphism
\[
 t=\frac{1-\zeta h}{h-1},\qquad
       (h,v)\longmapsto[t:1:-(t+1)/v].
 \tag{E80.17}
\]
Every denominator is invertible on the indicated
schemes. The equation for $v$ proves that this
point lies on the cubic; substitution proves both
compositions are the identity. The Kummer function
has orders $2,-1,-1$ at $0,-\zeta,\infty$,
respectively. Its three-valuation at $-\zeta$ proves
that it is not a cube. Removing these points and
$1$ makes the degree-three cover finite \'{e}tale and
connected. The filled points $x=0$ map to $h=\zeta^2$;
this is an unramified fibre of three points.

Apply the finite-cover comparison in
\cite[Expos\'e XII, Theorem 5.1 and Corollary 5.2]{E_SGA1}
to these actual morphisms. The topological punctured
curves give the free pro-three presentations just
stated. For completeness, the pro-three topology on
the index-three subgroup is its full one: intersect
the three conjugates of any normal subgroup of
three-power index in that subgroup. The intersection
is normal in the whole group, and the resulting
extension has three-power order. Thus the subgroup
comparison preserves every finite three-cover,
and its rank is $1+3(3-1)=7$. Filling the three
specified punctures gives the surjection
$N\to N_0$ with the closed normal closure of their
peripheral generators as its kernel.

On $S$ the first map fills $-\zeta$, and the second
fills $\infty$ before applying the displayed
M\"obius transformation. Their common quotient
fills both points. It is the pro-three group of
$\mathbb P^1\setminus\{0,1\}$, hence $\mathbb Z_3$.
The two maps to this quotient are exactly $\ell$:
the Kummer functions $h/(1-h)$ and
$\bar h/(1-\bar h)$ have ratio $-\zeta^{-1}$,
so their geometric characters agree. Given a pair
with equal common character, lift its first
coordinate to $H$ and correct its second coordinate
by the first kernel. This proves surjectivity onto
the fibre product, since the image of that kernel
is exactly the kernel of the second common-character
map. All these images and products are closed,
being images of compact groups.

Here is also a direct verification of the asserted
intersection kernel. Choose a free peripheral basis
of $H$ consisting of one surviving generator and
the two separately killed generators. Let $A,B$ be
the kernels of the retractions from the respective
free rank-two factors onto the surviving
$\mathbb Z_3$. The fibre product is
$(A\times B)\rtimes\mathbb Z_3$. In
$H/\overline{[K_\infty,K_{-\zeta}]}$, conjugation
of the first killed generator by the other kernel
does nothing. Its normal closure therefore comes
from $A$; the other normal closure comes from $B$.
The two commute, and the surviving generator acts
on both. This constructs the reverse continuous map
from $(A\times B)\rtimes\mathbb Z_3$. The two
compositions fix the three free generators, so are
the identity. Hence the intersection of the original
kernels is exactly their closed mutual commutator.

The character of the cover (E80.2) is
$X(f)+X(g)\bmod3$, since its actual defining root
is $(h\bar h)^{1/3}$. It vanishes on the intersection
kernel. Restricting the already proved fibre-product
surjection to this character kernel proves (E80.3).
For (E80.4), lift each prescribed $C_r$ coordinate
arbitrarily to $F$. Both required conditions are
abelian, so the pair of lifts lies in $\mathcal J$.
The stated kernel is then the literal inverse image
of $\gamma_{r+1}F$ under both maps, divided by
$\gamma_{r+1}N$.

\textbf{Roots, paths and the descended action.}
The chosen square roots satisfy
$(\lambda_n/\bar\lambda_n)^2=\zeta_n^2$.
At level zero their ratio is $-\zeta$; compatibility
and the odd cube power give $-\zeta_n$ at every
level. Their product is the positive root $3_n$.
The remaining multiplicative identities in (E80.5)
follow directly. Taking Kummer characters gives
\[
 \kappa_a+\kappa_{\bar a}=-3\kappa_3(3),\qquad
 \kappa_a-\kappa_{\bar a}=-\frac{\chi-1}{3},
 \tag{E80.18}
\]
which proves the two individual formulas with their
actual coherent roots.

Near $O$, the cubic equation gives
$y/x=-1-s^3/3+s^6/9+\cdots$ and
$v=s^2/3+\cdots$. The leading terms of $h,\bar h$
are $a s^3,\bar a s^3$. The product of the two
chosen cube-root leading coefficients is $1/3$.
The product cube root therefore has the same
leading term as $v$; their quotient is a cube root
of one with leading coefficient one, so it is one.
At every higher level the same argument identifies
the product roots with the corresponding coherent
roots of $v$. This proves the actual sheet and
tangential compatibility, including its constant.

E.78 transports the actual source marking so that
$\chi(\tau)=\chi(\tau')=7$ and
$\kappa_3(3)(\tau)=1$, $\kappa_3(3)(\tau')=-1$.
Substitution in (E80.5) gives the four parameters
in (E80.6). E.54's genuine operator, whose
class-three coefficient is $(\chi^2-1)/24$, gives
exactly E.78's polynomial $T_\mu$ at these values.
All Hall terms are therefore those of the actual
arithmetic action. For conjugate paths functoriality
is the exact identity
$\bar h_*\circ c_{\rm ar}=I\circ h_*$.
On the target the positive real interval is fixed
by complex conjugation and its peripheral circles
are reversed; thus $I(x)=x^{-1},I(y)=y^{-1}$.
The Hall formula is
\[
      I(A,B,n,r,s)=(-A,-B,n,-n-r,-n-s).
 \tag{E80.19}
\]
It proves both the stated involution and
$\jmath\mathcal T\jmath=\mathcal T'$ by direct
substitution. Each $T_\mu$ multiplies both abelian
characters by seven; $\jmath$ negates and swaps
them. Hence all preserve the actual joint image.
E.78's rational origin transport and E.34 identify
the based $N$ action with the same actual
$P\rtimes C_2$ action. This proves equivariance of
the quotient maps, rather than inferring a descent
from the first factor alone.

\textbf{Full local conjugacy and its inverse.}
The cohomological source comparison of E.74 and the
actual relative presentation of E.77 apply to this
pro-three target. Coprime complement conjugacy
normalizes the cocycle on $C_2$ to one. E.77 then
gives exactly (E80.7), with all residual conjugators
fixed by $\jmath$. A fixed pair has the form
$(w,Iw)$; its common-character condition says
$\ell(w)=-\ell(w)$, hence $\ell(w)=0$.
Its index-three condition is automatic. This proves
the final line of (E80.7).

Write $u_0=(u,Iu')$, $a_0=(a,Ia')$. The first
component of the relation in (E80.7) is exactly
$v_{-5/2}(u,a)=v_{1/2}(u',a')$, because of
$\jmath\mathcal T\jmath=\mathcal T'$. Thus
restriction gives E.78's actual source equation.
Membership of these two pairs is exactly (E80.9)
with the stronger equality $H_u=0$.

For an arbitrary E.78 gauge $w$, the abelian
coordinates transform by
\[
 \begin{aligned}
 \ell(u)+\ell(u')&\longmapsto H_u-12\ell(w),\\
 X(u)-X(u')&\longmapsto X(u)-X(u'),\\
 (X(a),Y(a),X(a'),Y(a'))&\longmapsto
                         (X(a),Y(a),X(a'),Y(a')).
 \end{aligned}
 \tag{E80.20}
\]
The existence of an integral normalizing gauge is
therefore exactly $H_u\in3\mathbb Z_3$, and the
displayed choice $w=x^{H_u/12}$ supplies it.
This proves necessity, sufficiency and the inverse
map in (E80.9)--(E80.10).

If two normalized joint representatives have the
same image in $H^1(k,C_r)$, an actual $w\in C_r$
relates their first quadruples. Both horizontal
$\ell$-sums are zero, so (E80.20) forces
$12\ell(w)=0$. The lattice $\mathbb Z_3$ is
torsion-free, hence $\ell(w)=0$. The pair $(w,Iw)$
is precisely a residual conjugator in (E80.7).
This proves injectivity on full geometric-conjugacy
classes, including the normalization which removes
the larger ambient target gauge. All spaces of
solutions are closed in finite products of compact
groups; quotienting a compact group action gives
a compact Hausdorff space. The displayed maps are
continuous, including division by twelve on its
closed image. The bijections and inverses are thus
homeomorphisms.

\textbf{Evaluated local coordinates and fibres.}
The map in (E80.11) is an integral parametrization
of the joint abelian lattice: $p=X(f)$,
$L=\ell(f)$, $m=(X(f)+X(g))/3$ recover its inverse.
The involution follows by negating and swapping
the two factors. Its fixed lattice is exactly the
$p$-axis. E.77's relation is identically zero on
this abelian target since $\mathcal T=7$.
The fixed gauge changes $u$ by minus six times
that axis and fixes $a$, proving the complete
coordinates and inverse in (E80.12).

Under (E80.10), the first-factor abelian values are
\[
 \begin{array}{ll}
 (A,B)=(p_a,p_a-L_a),&
 (A',B')=(p_a-3m_a,p_a-3m_a+L_a),\\
 (U,V)=(p_u,p_u-L_u),&
 (U',V')=(p_u-3m_u,p_u-3m_u+L_u).
 \end{array}
 \tag{E80.21}
\]
Their horizontal $\ell$-sum is already zero.
Substituting them into E.78's exact second-stage
defect gives
\[
 \begin{aligned}
 D_2={}&9L_am_a+3L_am_u-6L_ap_a-2L_ap_u+2L_a\\
       &-3L_um_a+2L_up_a-27m_a^2+18m_ap_a
                      +\tfrac32m_a-3p_a.
 \end{aligned}
 \tag{E80.22}
\]
Reduction modulo three is precisely the congruence
in (E80.12). E.78 proves that divisibility is
sufficient and gives its particular lift.

All the joint-image constraints depend only on
abelian coordinates. Consequently every E.78 lift
of the restricted normalized quadruple, at either
of the next two stages, still satisfies those
constraints and returns by (E80.10). Conversely
every joint lift restricts to that E.78 lift.
Injectivity already proved identifies the full
fibres, not only selected representatives. The
explicit central changes, stabilizer calculation
and inverse gauges in E.78 therefore give exactly
(E80.12)--(E80.13). In particular the coefficient
for a vertical third-degree change is $336$, and
the remaining horizontal diagonal gauge is $342$;
their three-adic valuations are one and two.
At the preceding step the corresponding coefficients
are $6$ and $48$, with valuations one and one.
The lower stabilizers vanish as in E.78: their
first nonzero Hall coordinate would be fixed by
$7^j$ for a positive integer $j$, impossible in
a torsion-free $\mathbb Z_3$ lattice.
Substitution of (E80.14) into the full E.78
polynomials gives $\Delta R=0$ and
$\Delta S=7(5-N-M)$, proving every assertion
about that descended family.

\textbf{Integral group-level lower central comparison.}
The following four actual elements of $\mathcal J_3$
will suffice:
\[
 \begin{gathered}
 g_1=(y^{-1},y^{-1}),\qquad
 g_2=(xy,x^{-1}y^{-1}),\\
 g_3=(1,x^3y^3),\qquad g_4=(c,1).
 \end{gathered}
 \tag{E80.23}
\]
The first three abelian columns are the three
columns of (E80.11). For any two elements their
second-degree commutator coordinates are the two
abelian determinants. On this integral lattice
these span the columns
\[
       \begin{pmatrix}-1&0&0\\1&-3&0\end{pmatrix},
       \qquad\text{Smith list }(1,3).
 \tag{E80.24}
\]
The resulting image is exactly
$\{(n,n'):n+n'\in3\mathbb Z_3\}$.

For the full class-three calculation use, in order,
$[g_4,g_1]$, $[g_4,g_2]$,
$[[g_1,g_2],g_1]$, $[[g_1,g_2],g_2]$.
They have zero first- and second-degree coordinates.
Their four central coordinate columns and their
integer inverse are
\[
 \begin{gathered}
 B=\begin{pmatrix}
 0&-1&0&-1\\1&-1&1&-1\\0&0&0&-1\\0&0&-1&-1
 \end{pmatrix},\qquad\det B=-1,\\
 B^{-1}=\begin{pmatrix}
 -1&1&-1&1\\-1&0&1&0\\0&0&1&-1\\0&0&-1&0
 \end{pmatrix}.
 \end{gathered}
 \tag{E80.25}
\]
Thus all four central directions actually belong
to $\gamma_2\mathcal J_3$. Combining these exact
commutators with (E80.24) proves both claims for
$\gamma_2$ in (E80.15). It also proves that
(E80.23) topologically generates $\mathcal J_3$:
the first three generators cover the abelian
pair lattice, $g_4$ and $[g_1,g_2]$ cover both
remaining $c$ coordinates, and (E80.25) covers
all central coordinates. Continuous
$\mathbb Z_3$ powers give the stated integral
coordinates, so this proves the group-level
generation and not only its rational graded image.

For $\gamma_3$ only the abelian coordinates of
the first input and the $c$ coordinates of the
second input enter. Bracketing the three columns
of (E80.11) with the two lattice basis vectors
$(1,-1),(0,3)$ gives the exact matrix
\[
 M=\begin{pmatrix}
 0&0&1&0&0&0\\
 -1&0&1&0&0&0\\
 0&0&1&-3&-3&9\\
 1&-3&1&-3&-3&9
 \end{pmatrix},\qquad
       \operatorname{Smith}(M)=(1,1,3,3).
 \tag{E80.26}
\]
The two rows
$(-1,0,1,0)$ and $(-2,1,0,1)$ kill its columns
modulo three and are independent there. Their
kernel has index nine, equal to the Smith index,
so it is exactly the image of $M$. This proves
the last line of (E80.15) integrally.
The stored certificate includes the complete
unimodular left and right Smith maps and both
inverses, as well as (E80.25).

It follows that the product-central kernel has
the two finite quotients in (E80.16). On the
second quotient the actual involution and
$\tau$ act by
\[
       [\jmath]=\begin{pmatrix}1&0\\2&-1\end{pmatrix}
                        \pmod3,
       \qquad [\tau]=I_2.
 \tag{E80.27}
\]
Indeed the central involution is minus the factor
swap, and $\tau$ is multiplication by $343$.
The abstract compact-module decompositions follow
by splitting their free $\mathbb Z_3$ quotients;
the exact maps and action are the preceding ones.
Finally surjectivity in (E80.4) implies
$\operatorname{im}(\gamma_3N)=\gamma_3\mathcal J_3$.
Since $\gamma_4N$ is killed, this proves the
stated index-nine image of the intrinsic
third-degree lattice.

The self-contained program
\texttt{a\_joint\_wild.py} verifies the eleven
scheme and tangent identities, the actual Hall
operators, all integral Smith maps and their
inverses. Separate finite Hall arithmetic
enumerates the complete joint images at class
one modulo $27$, class two modulo $9$, and class
three modulo $3$, of orders $6561$, $19683$,
and $6561$. It checks $1200$ arithmetic action
identities, $240$ exact restriction/gauge
normalizations and inverses, $62$ descended
second-stage sections, $13$ third-stage lifts,
and all $36$ displayed small family cases.
The accompanying JSON retains all matrices and
the complete finite-image hashes. These checks
verify the formulas used in the proofs above.
\end{proof}
